\section{Finite scalar acquisition of a complete patch}
\label{sec:acquisition}
We now work in the plane. Fix once and for all a nonnegative, even
$C^\infty$ probability kernel $\kappa$ supported in the unit disk, and
put $K=\|\nabla\kappa\|_1$. We may increase $K$ to be at least one.
Take, for example,
\begin{equation}\label{eq:design}
 t=\min(1,d_0/4),\qquad \sigma_0=\min(1,d_0/8),\qquad
 m=\left\lfloor\frac{D_0+2\sigma_0}{t}\right\rfloor+1,
 \qquad a=t(1,0).
\end{equation}
The chain is fixed by the numerical priors, not by an unknown boundary.
For $f_{x,\sigma}(q)=\sigma^{-2}\kappa((q-x)/\sigma)$ define
\[
 d_\sigma(x)=P_a(f_{x,\sigma})-P_{-a}(f_{x+a,\sigma}).
\]
Equation~\eqref{eq:balance} and Lemma~\ref{lem:zero} give
\begin{equation}\label{eq:smoothed-min}
 d_\sigma(x)=u_\sigma(x+a)-u_\sigma(x),\qquad
 u_\sigma(x)=-\min_{0\le k\le m}\sum_{j<k}d_\sigma(x+ja).
\end{equation}
There is no deconvolution, time derivative or noise amplification by
$\sigma^{-1}$ in this step.

\begin{proposition}[Stable recovery of occupation averages]
\label{prop:mean-stability}
At each point $x$ the $2m$ scalar means in \eqref{eq:smoothed-min}
determine $u_\sigma(x)$. If each mean has absolute error at most
$\epsilon$, the error is at most $2m\epsilon$, uniformly in $\sigma$.
Clipping the answer to $[0,1]$ cannot increase this error. For two
physical candidates whose corresponding means differ by at most $b$,
their occupation averages at $x$ differ by at most $2mb$.
\end{proposition}
\begin{proof}
Each balance difference has error at most $2\epsilon$, or at most
$2b$ in the comparison version. Apply the perturbation conclusion of
Theorem~\ref{thm:chain}. The true average lies in $[0,1]$.
\end{proof}

Each mean is estimated by repeated attempts with the prescribed input
law. The actual start position need not be recorded with the output;
the commanded distribution, its center and the direction are known.
The only per-attempt recorded variable is the bit. If the actual launch
law differs from its command by total variation at most $\zeta$, the
bias of every mean is at most $\zeta$. In particular a center error
$\ell$ changes this smooth input density in $L^1$ by at most
$K\ell/\sigma$. These observations describe input calibration costs;
they are not a claim that arbitrary directional or dynamical apparatus
errors are harmless.

\begin{lemma}[A finite grid controls the patch geometry]
\label{lem:grid}
Assume all physical bodies are strictly convex, have curvature at most
$\kappa_+$, and satisfy the common separation and diameter bounds.
Let $U$ be a square containing, with a collar of width $2D_0+2$, every
body to be compared. Take
$\sigma<\min(\sigma_0,(2\kappa_+)^{-1},d_0/20,1/4)$.
A grid of covering radius $\sigma/(32K)$ in $U$ contains
$O_U(\sigma^{-2})$ points. If the scalar means for two admissible
physical candidates differ at every grid command by at most
$b=1/(32m)$, then their complete bodies in the protected inner patch
can be matched with Hausdorff error at most $2\sigma$.
\end{lemma}
\begin{proof}
By Proposition~\ref{prop:mean-stability} the occupation averages differ
by at most $1/16$ at each grid point. Each occupation function has
Lipschitz constant at most $K/\sigma$. The chosen covering radius
therefore gives a uniform difference at most $1/8$ on $U$.

For a convex body $C$, let $C_{-\sigma}$ denote its erosion by the
closed radius-$\sigma$ disk. If $x\in C_{-\sigma}$, then
$u_\sigma(x)=1$, whereas if $\dist(x,\mathcal O')>\sigma$, the other
occupation average is zero. It follows that
\begin{equation}\label{eq:erosion-inclusion}
 \mathcal O_{-\sigma}\cap U\subset\mathcal O'+\sigma\overline B,
 \qquad
 \mathcal O'_{-\sigma}\cap U\subset\mathcal O+\sigma\overline B.
\end{equation}
The notation on the left means the union of the component erosions.
The curvature upper bound ensures an interior rolling radius
$r=1/\kappa_+$. One elementary justification uses the support function:
$p+p''\ge r$, so $p-r$ is a convex support function and
$C=C_0+r\overline B$. Hence
$C=(C_{-\sigma})+\sigma\overline B$ when $\sigma<r$.
Each erosion is nonempty and connected.

Distinct expanded target components have positive separation. A
connected erosion in \eqref{eq:erosion-inclusion} must therefore lie
in the radius-$\sigma$ enlargement of a single target body. Restoring
the disk shows $C\subset D+2\sigma\overline B$. Apply the other
inclusion to $D$ to obtain $D\subset E+2\sigma\overline B$ for a
single source body $E$. All these bodies and their erosions are in $U$
by the protective collar. Since $C\subset E+4\sigma\overline B$ and
$4\sigma<d_0$, necessarily $E=C$. Thus $d_H(C,D)\le2\sigma$.
The same argument in the reverse direction makes the matching bijective
for all complete bodies used in the protected patch. No incomplete
outer fragment is declared to be a physical component.
\end{proof}

For later use we record the smooth interpolation underlying the
$C^2$ conclusion. Here $C^j$ norms of support functions are on the
normal-angle circle, in a fixed laboratory frame.
\begin{lemma}[From Hausdorff error to smooth support error]
\label{lem:smoothing}
For convex bodies with bounded $C^6$ support norms, Hausdorff distance
at most $e\le1$ implies $C^2$ support difference at most $C e^{2/3}$.
The same estimate holds for centered support functions, with a changed
constant and their Steiner-center error bounded by $2e$.
\end{lemma}
\begin{proof}
Let $\varphi$ be a smooth even probability kernel of compact support
on the line, and set $\varphi_2(s)=\varphi(s/2)/2$ and
$J=(4\varphi-\varphi_2)/3$. Then $J$ has integral one and its moments
of degrees one, two and three vanish. Periodized convolution at scale
$h$ satisfies
\[
 \|J_h*(p-p')\|_{C^2}\le Ceh^{-2},\qquad
 \|J_h*p-p\|_{C^2}+\|J_h*p'-p'\|_{C^2}\le C h^4.
\]
The first inequality uses $\|p-p'\|_\infty=d_H(C,D)$ and the $L^1$
norms of the first two derivatives of the kernel. The second is Taylor's
formula for $p,p',p''$ through the cancelled orders, using the $C^6$
bound. Choose $h=e^{1/6}$. The Steiner point is
$\pi^{-1}\int_0^{2\pi}p(\theta)(\cos\theta,\sin\theta)\dd\theta$;
its change has norm at most $2e$. Subtracting this translation gives
the centered estimate.
\end{proof}

An admissible-candidate formulation will suffice below. It avoids
claiming an efficient reconstruction algorithm for all smooth shapes:
we select a physical candidate consistent with the finitely many scalar
means, then apply the explicit finite period tests to its recovered
patch. The statements concern both this selection and its quantified
error, not the execution time of searching an infinite-dimensional class.
